\section{TASK 3: MODEL TRAINING} \label{sec:task3}

The NER task is formulated as word-level MITRE tactic prediction after
subword tokenization. Each sub-token inherits its parent word's label for both
training and evaluation; only special tokens (\texttt{[CLS]}, \texttt{[SEP]},
padding) are excluded via the \texttt{-100} ignore index. The original
training split is further divided into \textbf{202 training sessions} and
\textbf{49 validation sessions} using multilabel stratification at session
level, so that rare tactics are represented more reliably during model
selection. The held-out test set contains \textbf{108 sessions}.

\subsection{Validation Strategy and Fine-Tuning Setup} \label{subsec:task3_setup}

For the full fine-tuning experiments, model selection is based on validation
macro F1 together with validation loss. Macro F1 is the main classification
criterion because the labelled dataset is imbalanced, while validation loss is
used to detect unstable convergence or late overfitting. All final numbers in
the tables are computed on the held-out test set only after selecting the
configuration from the validation histories.

\vspace{0.2em}

\noindent
\textbf{Pre-trained BERT.} The first model is \texttt{bert-base-uncased}
fine-tuned end-to-end for token classification. Despite the small labelled
dataset, the model reaches \textbf{84.76\% token accuracy}, \textbf{61.97\%
macro F1}, and \textbf{79.97\% average session fidelity} on the test set
(Table~\ref{tab:task3_performance}). This is a useful result, but the gap
between accuracy and macro F1 shows that the model is mostly strong on frequent
classes and less reliable on minority tactics.

Figure~\ref{fig:task3_bert_lr} shows the validation behaviour for the BERT
learning-rate search. The smallest rate is stable but slow, while larger rates
improve faster and then become more irregular. The selected checkpoint is
therefore not simply the last epoch, but the epoch that gives the best
validation compromise before the model starts fitting the small training split
too closely.

\begin{figure}[H]
    \centering
    \includegraphics[width=0.88\linewidth]{\imagespath Task3/bert_pretrained_lr_comparison.png}
    \caption{Validation loss and macro F1 curves for pre-trained BERT across
    learning rates.}
    \label{fig:task3_bert_lr}
\end{figure}

\noindent
\textbf{Naked BERT baseline.} To isolate the contribution of pre-training, the
same BERT architecture is trained from random initialization. Its macro F1 drops
to \textbf{44.17\%}, confirming that the task is not solved equally well by any
large architecture when only 251 labelled sessions are available. Pre-training
therefore provides a strong prior, even though it was learned mostly on natural
language rather than shell logs.

\vspace{0.2em}

\noindent
\textbf{UnixCoder.} The code-oriented \texttt{microsoft/unixcoder-base} model is
the best full fine-tuned model. It reaches \textbf{87.10\% accuracy},
\textbf{73.50\% macro F1}, and \textbf{86.29\% fidelity}, improving over
pre-trained BERT by \textbf{11.53 macro-F1 points}. This supports the hypothesis
that a model pre-trained on code-like data is better suited to SSH sessions,
which contain commands, paths, flags, separators, and encoded strings rather
than natural prose.

The UnixCoder validation curves in Figure~\ref{fig:task3_unixcoder_lr} show a
clearer separation between useful and less useful learning rates. The selected
configuration reaches the best validation macro F1 without requiring the very
late epochs, which is important because the training set has only a few hundred
sessions. This behaviour supports using UnixCoder as the backbone for the
subsequent frozen-layer experiments and for inference.

\begin{figure}[H]
    \centering
    \includegraphics[width=0.88\linewidth]{\imagespath Task3/unixcoder_lr_comparison.png}
    \caption{Validation loss and macro F1 curves for UnixCoder across learning
    rates.}
    \label{fig:task3_unixcoder_lr}
\end{figure}

\begin{table}[H]
    \centering
    \small
    \begin{tabular}{lrrrrrr}
        \toprule
        \textbf{Model} & \textbf{Acc.} & \textbf{Prec.} & \textbf{Rec.} &
        \textbf{Macro F1} & \textbf{Fidelity} & \textbf{Val F1} \\
        \midrule
        BERT pre-trained & 84.76 & 75.68 & 56.51 & 61.97 & 79.97 & 73.23 \\
        Naked BERT & 73.95 & 67.68 & 41.08 & 44.17 & 75.06 & 56.53 \\
        UnixCoder & \textbf{87.10} & \textbf{86.63} & \textbf{68.80} &
        \textbf{73.50} & \textbf{86.29} & \textbf{79.88} \\
        UnixCoder Last2+Head & 83.89 & 75.36 & 61.11 & 64.85 & 83.05 & 72.24 \\
        UnixCoder Head Only & 76.98 & 68.13 & 48.50 & 52.57 & 76.34 & 59.41 \\
        \bottomrule
    \end{tabular}
    \vspace{0.1cm}
    \caption{Test performance and validation score for the trained NER models.}
    \label{tab:task3_performance}
\end{table}

\vspace{-0.1cm}

\begin{figure}[H]
    \centering
    \includegraphics[width=0.78\linewidth]{\imagespath Task3/model_comparison.png}
    \caption{Comparison of token accuracy, macro F1, and session fidelity across
    the trained models.}
    \label{fig:task3_model_comparison}
\end{figure}

\subsection{Per-Class Behavior} \label{subsec:task3_per_class}

The per-class F1 plots in Figure~\ref{fig:task3_per_class} make the imbalance
effect visible. Frequent tactics such as Discovery and Execution are learned
more consistently, while classes with few examples remain less stable. This is
the main difficulty of the task: high token accuracy can coexist with weaker
balanced performance because most words belong to a small number of tactics.
UnixCoder improves this behavior, but it does not remove the scarcity problem
for the rarest classes.

\begin{figure}[H]
    \centering
    \begin{subfigure}[b]{0.48\linewidth}
        \centering
        \includegraphics[width=\linewidth]{\imagespath Task3/bert_pretrained_per_class_f1.png}
        % \caption{BERT pre-trained.}
    \end{subfigure}
    \hfill
    \begin{subfigure}[b]{0.48\linewidth}
        \centering
        \includegraphics[width=\linewidth]{\imagespath Task3/unixcoder_per_class_f1.png}
        % \caption{UnixCoder.}
    \end{subfigure}
    \caption{Per-class F1 scores for (a) BERT and (b) UnixCoder fine-tuned pre-trained models.}
    \label{fig:task3_per_class}
\end{figure}

\noindent
The normalized confusion matrix in Figure~\ref{fig:task3_unixcoder_cm} provides
a more detailed view of the best model. The dominant classes are recognized
reliably, which explains the high token accuracy and session fidelity. The
remaining errors are concentrated in semantically close or low-support regions:
some command chains with mixed behaviour are confused between Discovery and
Execution, while the rarest tactics remain harder to separate because the model
sees only a small number of labelled examples during fine-tuning.

\begin{figure}[H]
    \centering
    \includegraphics[width=0.78\linewidth]{\imagespath Task3/unixcoder_confusion_matrix_normalized.png}
    \caption{Normalized confusion matrix for the fully fine-tuned UnixCoder
    model on the test set.}
    \label{fig:task3_unixcoder_cm}
\end{figure}

\noindent
This error profile is coherent with the structure of the dataset. Many SSH
sessions are long pipelines where a word can support more than one adversarial
intent depending on the downstream command. The model is therefore not only
learning command names, but also trying to infer the role of each word in a
local execution context. The confusion matrix confirms that UnixCoder is the
most reliable candidate, but it also motivates keeping per-class metrics in the
evaluation rather than relying only on aggregate accuracy.

\subsection{Frozen-Layer Experiments} \label{subsec:task3_frozen}

The best full model is selected using validation macro F1 only; the selected
backbone is UnixCoder. Two reduced fine-tuning strategies are then tested:
training only the last two encoder layers plus the classification head, and
training only the classification head. Table~\ref{tab:task3_training_cost}
summarizes the computational trade-off.

Freezing most of the encoder reduces trainable parameters from
\textbf{125.3M} to \textbf{14.2M} and nearly halves training time, from
\textbf{686.5 s} to \textbf{349.0 s}. However, macro F1 decreases from
\textbf{73.50\%} to \textbf{64.85\%}. Training only the head is faster
(\textbf{285.9 s}) and uses only \textbf{5,383 trainable parameters}, but the
macro F1 falls to \textbf{52.57\%}. Therefore, the pre-trained representation is
useful, but adapting part of the encoder is still necessary for this domain.
The frozen settings also require larger learning rates than full fine-tuning,
especially for the randomly initialized classification head.

\vspace{0.2cm}

\begin{table}[H]
    \centering
    \small
    \begin{tabular}{lrrrrr}
        \toprule
        \textbf{Model} & \textbf{Base LR} & \textbf{Head LR} &
        \textbf{Trainable} & \textbf{Epoch} & \textbf{Time (s)} \\
        \midrule
        BERT pre-trained & 0.00001 & 0.00005 & 108,897,031 & 20 & 761.0 \\
        Naked BERT & 0.00005 & 0.00005 & 108,897,031 & 13 & 764.0 \\
        UnixCoder & 0.00002 & 0.00010 & 125,344,519 & 11 & 686.5 \\
        UnixCoder Last2+Head & 0.00005 & 0.00025 & 14,181,127 & 16 & 349.0 \\
        UnixCoder Head Only & 0.00100 & 0.00100 & 5,383 & 21 & 285.9 \\
        \bottomrule
    \end{tabular}
    \vspace{0.1cm}
    \caption{Training configuration and cost for the Task~3 models.}
    \label{tab:task3_training_cost}
\end{table}
